\begin{table}

\caption{\label{tab:it-hbs-latent-income-coefficients}Latent income-score model, Italy}
\centering
\begin{tabular}[t]{lcc}
\toprule
Variable & 2015 & 2020\\
\midrule
Absolute-poverty indicator & 0.142 & 0.001\\
Age of reference person & 0.139 & 0.007\\
Consumption level & 0.157 & -0.017\\
\addlinespace
Economic resources & 0.144 & 0.017\\
Education of reference person & 0.151 & 0.010\\
\addlinespace
Housing surface per person & 0.149 & -0.006\\
Housing tenure & 0.145 & 0.002\\
Labour-market status & 0.137 & -0.000\\
\addlinespace
Occupation & 0.143 & 0.001\\
Subjective economic situation & 0.145 & -0.002\\
\bottomrule
\end{tabular}
\end{table}
\begin{flushleft}
\footnotesize{Notes. The table reports coefficients of the latent socio-economic score used to construct probabilistic income groups in the 2015 and 2020 Istat HBS public-use files. The model excludes household-size, age-squared, rooms-per-person variables and region or macro-region fixed effects. Continuous variables are standardised. The coefficients should be read as score weights used for basket construction rather than causal estimates. Sources: Istat HBS public-use files, Eurostat HBS, authors' calculations.}
\end{flushleft}
